\documentclass[11pt,a4paper]{article}
\usepackage[margin=2.4cm]{geometry}
\usepackage{amsmath,amssymb,booktabs,hyperref,listings,xcolor}
\lstset{basicstyle=\ttfamily\footnotesize,breaklines=true,frame=single,
        backgroundcolor=\color{gray!8}}
\title{\textbf{Plan --- Phase 2: reverse-mode (adjoint) gradients for rotor chains}\\
       \texttt{hymeko\_clifford::autodiff}}
\author{Aiko Seto}
\date{Created-at: 2026-08-16 20:55 CEST}

\begin{document}
\maketitle

\section*{ETA}
$\approx$50 min: 20 min module + contracts, 15 min tests (adjoint vs finite
difference vs closed form, operation counts), 10 min benchmark + measurement
against the recorded $O(k^2)$ baseline, 5 min report. Basis: the mathematics is a
two-sweep recurrence with a known closed form; the correctness harness
(\texttt{tests/scenario\_derivative.rs}) and the cost baseline
(\texttt{clifford/derivative/*}) both already exist, measured 2026-08-16.

\section{Scope}
\texttt{hymeko\_clifford} is documented as the ``Clifford-algebra autograd
backend for G-SPHF''; Phase 1 shipped no gradients. Measured consequence
(\texttt{reports/2026-08-16-nagare-cpu-memory-benchmark.md} \S6h): a
central-difference gradient over a $k$-angle rotor chain costs $2k$ forward
evaluations, i.e.\ $O(k^2)$ --- $63.9\times$ the forward pass at $k=32$.

This plan adds the reverse sweep, which is $O(k)$.

\textbf{Not in scope:} a general autograd tape over arbitrary multivector
expressions. The rotor chain is the shape that is measured, benchmarked, and
used; a tape is a larger decision and should follow evidence from this one.

\section{Mathematics}
With $x_0 = v$, $x_i = R_i x_{i-1} \tilde R_i$, $w = x_k$, and rotors
$R_i = \exp(-\tfrac{\theta_i}{2} B_i)$ about unit axes $\omega_i$:
\[
  \frac{\partial x_i}{\partial \theta_i} = \omega_i \times x_i,
  \qquad
  \frac{\partial x_i}{\partial x_{i-1}} = R_i(\cdot)\tilde R_i .
\]
Seeding the adjoint with $\bar a_k = \partial L/\partial w$, the reverse
recurrence is
\[
  \frac{\partial L}{\partial \theta_i} = \bar a_i \cdot (\omega_i \times x_i),
  \qquad
  \bar a_{i-1} = \tilde R_i\, \bar a_i\, R_i ,
\]
and $\bar a_0 = \partial L/\partial v$ falls out for free. One forward sweep
storing $x_0 \dots x_k$, one backward sweep: $2k$ rotor applications and $k$ cross
products, \textbf{$O(k)$ in both time and memory}, against $2k^2$ rotor
applications for the central difference (already asserted by counter in
\texttt{gradient\_of\_a\_chain\_costs\_two\_evaluations\_per\_angle}).

\section{Affected files}
\begin{itemize}
  \item \texttt{src/autodiff.rs} --- new module: \texttt{RotorChain},
        \texttt{Adjoint}, \texttt{ChainError}.
  \item \texttt{src/lib.rs} --- declare and re-export; the Phase-1 doc comment
        updated so it no longer claims gradients are absent.
  \item \texttt{tests/scenario\_adjoint.rs} --- new scenario tests.
  \item \texttt{../../nagare/benches/clifford\_algebra\_bench.rs} --- adjoint arm
        beside the existing finite-difference sweep, same shapes.
  \item Mirror every source change to
        \texttt{nagare/vendor/hymeko\_clifford/}, which is byte-identical today
        and must stay so.
\end{itemize}

\section{CORE.YAML items touched}
\textbf{None} --- \texttt{CORE.YAML} does not exist in this workspace (flagged
2026-08-16). \textbf{No dependency added}: the module is \texttt{std}-only, like
the rest of the crate.

\section{Interface}
\begin{lstlisting}[language=Rust]
pub struct RotorChain { /* axes: Vec<[f64; 3]> */ }
pub struct Adjoint { pub value: [f64; 3],
                     pub grad_angles: Vec<f64>,
                     pub grad_input: [f64; 3] }
pub enum ChainError { NonUnitAxis { index, norm }, AngleCount { expected, got } }

impl RotorChain {
    pub fn new(axes: Vec<[f64; 3]>) -> Result<Self, ChainError>;
    pub fn apply(&self, angles: &[f64], v: [f64; 3]) -> Result<[f64; 3], ChainError>;
    pub fn adjoint(&self, angles: &[f64], v: [f64; 3],
                   grad_out: [f64; 3]) -> Result<Adjoint, ChainError>;
    pub fn finite_difference_gradient(&self, angles: &[f64], v: [f64; 3],
                                      grad_out: [f64; 3], h: f64)
        -> Result<Vec<f64>, ChainError>;
}
\end{lstlisting}

\texttt{Result} rather than panics: axis normalisation and angle-count are
caller errors, and \S6.4 forbids swallowing them. The finite difference ships
\emph{in the library} deliberately --- it is the reference a caller uses to
gradient-check its own objective, and keeping it beside the adjoint means the
two cannot drift apart unnoticed.

\section{Test strategy}
\begin{description}
  \item[Unit] adjoint vs the closed form $\omega \times w$ at $k=1$; error type
        surfaces a non-unit axis and an angle-count mismatch.
  \item[Scenario] adjoint vs central difference over $k \in \{1,2,4,8,16\}$,
        random axes and seeds; \texttt{grad\_input} checked against a finite
        difference in $v$; gradient of a closed loop is zero where the loop is
        the identity.
  \item[Cost] operation counters assert $2k$ rotor applications for the adjoint
        against $2k^2$ for the difference --- the $O(k)$ claim machine-checked,
        not left to the clock.
  \item[Regression] the adjoint did not exist, so every test is new; the
        finite-difference reference is the oracle.
\end{description}

\section{Performance budget}
\begin{center}
\begin{tabular}{lrr}
\toprule
Metric ($k=32$, measured 2026-08-16) & Finite difference & Target (adjoint) \\
\midrule
gradient / forward ratio & $63.9\times$ & $\le 4\times$ \\
absolute gradient cost   & 51.42\,\textmu s & $\le 4$\,\textmu s \\
growth per doubling of $k$ & $\approx 4\times$ & $\approx 2\times$ \\
\bottomrule
\end{tabular}
\end{center}
Targets follow from the operation counts, not from optimism: the adjoint does
$2k$ rotor applications where the difference does $2k^2$. If the measurement
lands outside, the implementation is wrong or the constant is dominated by
allocation --- both are findings, not excuses.

\section{Rollback}
Purely additive: a new module, new tests, a new benchmark arm. Deleting
\texttt{src/autodiff.rs} and its \texttt{pub mod} line restores the crate exactly.
No existing signature changes, no data or fixture is mutated.

\section{Risk anticipation}
\begin{itemize}
  \item \textbf{Worst-case input}: long chains ($k$ up to $32$ benchmarked, no
        upper bound in the API), angles near $\pi$ where a Cayley
        parameterisation would diverge --- avoided by parameterising the rotor
        directly as $\cos(\theta/2) + \sin(\theta/2) B$ rather than through
        \texttt{cayley\_to\_unit\_quat}.
  \item \textbf{Precision}: the reference finite difference in \texttt{f64} is
        second-order accurate (verified: halving $h$ cuts error $4.00\times$).
        The pre-existing benchmark differences the \texttt{f32} quaternion path
        with $h=10^{-5}$, which sits near the cancellation limit for
        \texttt{f32}; that is a flaw in the \emph{baseline benchmark}, recorded
        in the report, and the reason this module and its comparison are
        \texttt{f64} throughout.
  \item \textbf{Contract preservation}: the crate is \texttt{\#![forbid(unsafe\_code)]}
        and \texttt{\#![warn(missing\_docs)]}; the new module inherits both, so
        every public item carries preconditions and postconditions.
  \item \textbf{Duplication}: two copies of the crate exist. Both are patched and
        re-checked with \texttt{diff -rq}, as for the \texttt{GradRequest} fix.
\end{itemize}

\section{Known deviation}
\textbf{\texttt{plan.pdf} cannot be produced}: no TeX toolchain on this host, and
installing one is a \S1 dependency change. Per \S11 this is a halt condition; it
is recorded and reported rather than skipped. The \texttt{.tex}, \texttt{.tikz}
and \texttt{.mmd} sources are complete.

\end{document}
